% ECONOMICS_PROBLEM: {"id": "Q58-61", "title": "Economics and governance of solar geoengineering and carbon removal", "jel_code": "Q58", "source_id": "OP-0120"}
% !TeX program = xelatex
\documentclass[11pt,a4paper]{article}
\usepackage[margin=23mm]{geometry}
\usepackage{fontspec}
\setmainfont{Latin Modern Roman}
\usepackage{amsmath,amssymb,mathtools}
\usepackage{xcolor,enumitem,booktabs,longtable,array,xurl,hyperref}
\definecolor{muted}{HTML}{53636C}
\hypersetup{colorlinks=true,urlcolor=blue,linkcolor=blue}
\setlength{\parindent}{0pt}
\setlength{\parskip}{5pt}
\newcommand{\R}{\mathbb R}
\newcommand{\E}{\mathbb E}
\newcommand{\N}{\mathbb N}
\newcommand{\ind}{\mathbf 1}
\newcommand{\Var}{\operatorname{Var}}
\newcommand{\Cov}{\operatorname{Cov}}
\newcommand{\supp}{\operatorname{supp}}
\DeclareMathOperator*{\argmin}{arg\,min}
\DeclareMathOperator*{\argmax}{arg\,max}
\newcommand{\entrymeta}[3]{{\sffamily\small\color{muted}\textbf{\texttt{#1}}\quad|\quad #2\par #3\par}}
\newcommand{\Problemref}[1]{\texttt{#1}}
\newcommand{\JELref}[2]{\texttt{#2} (JEL \texttt{#1})}
\begin{document}
\section*{Economics and governance of solar geoengineering and carbon removal}
\textbf{Problem ID:} \texttt{Q58-61}\quad \textbf{JEL:} \texttt{Q58}\par
% BEGIN ECONOMICS STATEMENT
\label{problem:OP-0120}
\label{atlas:prob:C10}

\noindent\textbf{Original identifier:} C10 (atlas item 120). \textbf{Field:} Environmental and climate economics. \textbf{Classification:} editorial research family with established restricted solutions; current open status not certified.

\subsection*{Definitions and assumptions}

Countries $i=1,\ldots,I$ choose mitigation $\mu_{i,t}$, carbon removal $r_{i,t}$ and solar intervention $g_{i,t}$ over a declared horizon. A carbon-cycle model tracks atmospheric and reservoir stocks, with emissions and removal in tonnes of $\mathrm{CO}_2$ per period; lifecycle emissions enter emissions explicitly and storage leakage returns carbon through stated transition equations. A distinct solar-intervention state determines radiative forcing in $\mathrm W/\mathrm m^2$ and can decay rapidly after deployment stops. Country payoff $W_i$ is expected discounted utility of consumption and climate/ecological outcomes, with technology costs, termination risk and regional heterogeneity. Each country's discounting, uncertainty law and information set are specified. Require every national payoff $W_i(e,m)$ to be a finite real number and each relevant equilibrium set to be nonempty.

\subsection*{Formal research question}

For this dynamic game, characterize equilibrium deployment, welfare effects and implementable governance mechanisms. Let $R$ specify consent, voting, transfers, liability, monitoring and enforceable sanctions, and let $\mathcal E(R,m)$ denote its stated equilibrium concept in model $m$. Seek rules satisfying participation relative to a declared outside-option equilibrium, incentive compatibility where private information is present, budget feasibility and technological safety constraints. Compare the set of national welfare vectors
\[
\mathcal W(R)=
\{(W_1(e,m),\ldots,W_I(e,m)):
m\in\mathcal M,\ e\in\mathcal E(R,m)\}
\]
across rules, where $\mathcal M$ encodes admissible physical and strategic uncertainty. Determine conditions for unilateral solar deployment, removal underprovision, mitigation substitution and negotiated improvements. Compare rules within matched models: $R$ weakly robust-Pareto-dominates $R^0$ if, for every $m\in\mathcal M$, every $e\in\mathcal E(R,m)$, every $e^0\in\mathcal E(R^0,m)$ and every country $i$, $W_i(e,m)\ge W_i(e^0,m)$. A strict robust improvement additionally requires some strict country inequality for each such matched equilibrium pair. Alternatively impose a stated selection $e_R(m)$ and compare selected pairs within each model. Unpaired welfare vectors from different models are not ranked as if their utility scales were interchangeable; aggregate national welfare only with explicit weights.

\subsection*{Interpretation and scope}

Solar radiation modification changes forcing without necessarily reducing atmospheric carbon or ocean acidification; carbon removal changes carbon stocks and faces permanence, additionality and lifecycle accounting. They therefore require distinct state dynamics rather than one generic deployment variable. A low engineering cost can create a free-driver externality, but does not make planetary damages small. A globally optimal control path assumes welfare weights and a decision-maker; Nash or bargaining predictions address a different institution. Monitoring ability does not imply sanction enforceability. An intervention's temperature target may leave precipitation, regional impacts or abrupt termination highly uncertain. Equilibrium selection and incomplete knowledge can require set-valued welfare rankings.

\subsection*{Source audit and status}

[S1] treats geoengineering economics, [S2] develops a free-driver voting architecture, and [S3] reviews geoengineering history and prospects. These verified sources support solar-governance motivation, with limited coverage of contemporary carbon-removal accounting and regulation. The combined two-technology dynamic-game question is editorial rather than a conjecture stated in those works. No conclusion about its exact October 2026 open status is inferred solely from these old references.

\subsection*{References}

\begin{enumerate}[label={[\textnormal{S}\arabic*]},leftmargin=*]

\item Scott Barrett (2008). \emph{The Incredible Economics of Geoengineering}. Environmental and Resource Economics 39(1), 45–54. \url{https://doi.org/10.1007/s10640-007-9174-8}. \textbf{Locator:} Publisher or institutional abstract and bibliographic record. \textbf{Support and limits:} Geoengineering deployment incentives and international governance; not comprehensive carbon-removal accounting.

\item Martin L. Weitzman (2015). \emph{A Voting Architecture for the Governance of Free-Driver Externalities, with Application to Geoengineering}. Scandinavian Journal of Economics 117(4), 1049–1068. \url{https://doi.org/10.1111/sjoe.12120}. \textbf{Locator:} Publisher or institutional abstract and bibliographic record. \textbf{Support and limits:} Free-driver externality and voting architecture in a specified model; does not guarantee real-world enforceability.

\item David W. Keith (2000). \emph{Geoengineering the Climate: History and Prospect}. Annual Review of Energy and the Environment 25, 245–284. \url{https://doi.org/10.1146/annurev.energy.25.1.245}. \textbf{Locator:} Publisher or institutional abstract and bibliographic record. \textbf{Support and limits:} Historical and conceptual review; publisher now labels the series Environment and Resources, while the original volume was Energy and the Environment.

\end{enumerate}

\subsection*{Common conventions: Source mathematical conventions}

Unless an entry states otherwise, agent and firm sets are finite. State and action spaces are measurable, strategies and outcome maps are measurable, probability laws are specified, and expectations exist for the targets under discussion. Infinite-horizon discount factors lie in $(0,1)$ and discounted objectives are finite. Norms, tolerances and complexity input sizes must be specified for computational comparisons. Constants or primitives that appear locally are inputs, not empirical facts asserted by this catalogue.

\textbf{Identification.} A statistical model $m\in\mathcal M$ implies an observable law $P_m$ and target $\theta(m)$. At law $P$, its identified set is
\[
\Theta(P;\mathcal M)=\{\theta(m):m\in\mathcal M,\ P_m=P\}.
\]
Point identification holds when this set is a singleton, or equivalently when $P_{m_1}=P_{m_2}$ implies $\theta(m_1)=\theta(m_2)$ for all compatible models. Sharp bounds mean bounds attained, or approached when the boundary is not attained, by compatible models. Writing this set is a definition, not a solution: the substantive task is to establish informative restrictions and derive or estimate the set. An unrestricted-confounding model may give only trivial bounds.

\textbf{Inference.} Identification, estimability and valid uncertainty quantification are separate questions. Fix $\alpha\in(0,1)$. A confidence procedure $C_n$ over a declared law class $\mathcal P$ has asymptotic uniform coverage at level $1-\alpha$ when
\[
\liminf_{n\to\infty}\inf_{P\in\mathcal P}\Pr_P\{\theta(P)\in C_n\}\geq1-\alpha.
\]
For set-identified targets the coverage event must be defined separately, for example coverage of the full identified set. If an entry uses identification-indexed models instead of uniquely indexed laws, the same coverage convention is applied over those models. Pointwise limits do not establish uniform validity, and asymptotic validity does not guarantee finite-sample precision. Sample sizes, dependence structures and weak-identification sequences are part of each application.

\textbf{Counterfactuals.} $Y_i(a)$ is a potential outcome under a specified intervention, including its timing, implementation and versions. It is not $Y_i$ conditional on voluntarily choosing $a$. A full intervention vector permits interference; shorthand scalar treatment notation does not silently impose no interference. Assignment, selection, attrition, anticipation and measurement must be specified. A claim about an instrument's local effect is not automatically a population policy effect.

\textbf{Economic equilibrium.} A counterfactual model includes best responses, constraints, feasibility, market clearing, state transitions and expectations consistent with the stated information. When several equilibria exist, either a declared selector is an input or the target is set-valued. Comparisons across policy regimes must use the stated selection convention. Reduced-form relationships are not presumed invariant after an equilibrium-changing intervention.

\textbf{Optimization and welfare.} Optimization is over the stated feasible set. If attainment is not established, $\sup$ or $\inf$ and $\varepsilon$-optimal choices replace an asserted maximizer or minimizer. For example, with value $V=\sup_{a\in\mathcal A}W(a)<\infty$, an $\varepsilon$-optimal set is $\{a\in\mathcal A:W(a)\geq V-\varepsilon\}$ for $\varepsilon>0$. Benefits, costs, risk and health or environmental outcomes can be aggregated only using declared units and conversion weights. Interpersonal utility weights, the treatment of fiscal transfers and foreign profits, discounting, and the behavioral welfare criterion are explicit inputs. A comparative-static derivative additionally requires existence and appropriate differentiability of the selected counterfactual.

\textbf{Sources.} Local labels [S1], [S2], and so forth restart in every question. Locators identify the relevant abstract, model, section or result at the level actually checked; a bibliographic or abstract-level verification is not represented as inspection of an entire proof. Working papers are distinguished from later publications and changing versions. Source summaries are paraphrased. All URL checks and editorial formulations refer to the audit date stated above.
% END ECONOMICS STATEMENT
\end{document}
